% All curves below share the initial conditions used in the adjacent solution.
\newcommand{\rTwoDataPath}{/Users/dongxingxi/Documents/物理竞赛一轮资料/sample/xes-original/thirty-eighth-section/assets/trajectory-data}
\newenvironment{rTwoFigure}{\par\noindent\begin{minipage}{\linewidth}\centering}{\end{minipage}\par\medskip}
\newcommand{\cyclotronExtractionDiagram}{\begin{rTwoFigure}
\begin{tikzpicture}[x=.9cm,y=.9cm,>=Latex]
 \begin{scope}[xshift=-3.4cm]
  \draw[thick](-.2,.7)--(-2,.7)--(-2,-.7)--(-.2,-.7);
  \draw[thick](.2,.7)--(2,.7)--(2,-.7)--(.2,-.7);
  \foreach \x in {-1.2,0,1.2}{\draw[vector](\x,1.3)--(\x,-1.3);}
  \draw[DeepBlue,densely dashed](-1.9,0)--(1.9,0);
  \node at(0,1.7){D 形电极的侧剖面};
  \node[align=center,text width=4.8cm,font=\footnotesize] at(0,-1.9){束流在中空腔内\\$\boldsymbol B$垂直轨道平面};
 \end{scope}
 \begin{scope}[xshift=3.2cm]
  \draw[guide](-2,0)..controls(-.6,0)and(0,.8)..(1,1.3);
  \draw[DeepBlue,thick,->](-2,0)..controls(-.5,0)and(.2,-.3)..(2,-.3);
  \draw[thick](-.9,-.8)--(.3,-.8);
  \node[below=3pt,font=\footnotesize]at(-.3,-.8){负偏转板};
  \draw[thin](1,-.55)--(1,.25);
  \draw[thin](1,.25)--(1,1.6)node[above]{薄铝窗};
  \draw[very thick](2,-.75)--(2,.2)node[above]{靶};
  \node[above] at(-1.1,.4){虚线：未引出时};
  \node[align=center,text width=4.8cm,font=\footnotesize] at(0,-1.9){实线：正离子被引向窗口\\再轰击靶};
 \end{scope}
\end{tikzpicture}
\end{rTwoFigure}}
\renewcommand{\magneticMomentDiagram}{\begin{rTwoFigure}
\begin{tikzpicture}[scale=.84,>=Latex]
\begin{scope}[xshift=-3cm]
 \draw[curve](0,0)circle(1.25);\fill(0,1.25)circle(2pt)node[above]{$+q$};
 \draw[vector](0,1.25)--(1,1.25)node[right]{$v_\perp$};
 \node at(-.2,.35){$\fieldOut\ B$};\node at(.2,-.35){$\fieldInto\ \mu$};
 \node at(0,-1.7){$\mu=E_\perp/B$};
\end{scope}
\begin{scope}[xshift=.8cm]
 \foreach \a in {-1.4,-.9,0,.9,1.4}{
  \draw[TextGray,->,domain=0:6,samples=60,variable=\t]plot(\t,{\a/sqrt(1+.15*\t)});
 }
 \draw[DeepBlue,thick]plot file {\rTwoDataPath/mirror.dat};
 \node[above]at(3,1.6){向右磁场逐渐增强};
 \node[align=center]at(3,-2){回旋半径减小，$v_\perp$增大\\$v_\parallel$减小；图为缓变近似投影};
\end{scope}
\end{tikzpicture}\end{rTwoFigure}}
\renewcommand{\gyroHelixDiagram}{\begin{rTwoFigure}
\begin{tikzpicture}[scale=.76,>=Latex]
\begin{scope}[xshift=-3cm]
 \draw[curve] (0,0) circle(1.4);
 \node at(0,0){$\fieldOut\ B$};
 \foreach \a in {90,210,330}{
  \fill[DeepBlue] (\a:1.4) circle(2pt);
  \draw[vector] (\a:1.4)--++({.65*sin(\a)},{-.65*cos(\a)});
  \draw[vector,TextGray] (\a:1.4)--(\a:.8);
 }
 \node[above] at(0,1.55){$q>0$};
 \node[below,align=center] at(0,-1.7){切向箭头：速度\\向心箭头：磁力};
\end{scope}
\begin{scope}[xshift=1cm]
 \draw[axis] (-.4,0)--(6.6,0) node[right]{$\boldsymbol B$};
 \draw[DeepBlue,thick,domain=0:1440,samples=450,variable=\t]
 plot({.0042*\t+.70*(cos(\t)-1)},{.78*sin(\t)});
 \fill[DeepBlue] (0,0) circle(2pt) node[below left]{$P$};
 \draw[vector] (0,0)--(1,0);
 \draw[thin] (.8,0)--(1.1,-.95)--(2.3,-.95) node[right]{$v_\parallel$};
 \draw[vector] (0,0)--(0,1.12) node[above]{$v_\perp$};
 \draw[<->] (0,-1.2)--(1.512,-1.2) node[midway,below]{$h$};
 \node at(3,-2){螺旋线的斜投影};
\end{scope}
\end{tikzpicture}
\end{rTwoFigure}}

\renewcommand{\driftDiagram}{\begin{rTwoFigure}
\begin{tikzpicture}[x=.70cm,y=.65cm,>=Latex]
 \foreach \row/\sgn/\name in {0/1/{质子},1/-1/{电子}}{
  \begin{scope}[yshift={-3.4*\row cm}]
   \draw[DeepBlue,thick,domain=0:720,samples=240,variable=\t]
    plot({.50*(pi*\t/180-sin(\t))},{\sgn*.50*(1-cos(\t))});
   \draw[->] (1.8,-1.4)--(3.5,-1.4);\node at(2.6,-1.8){$\boldsymbol v_E$};
   \node[anchor=east]at(-.3,0){\name};
  \end{scope}}
 \node at(3,2){$\boldsymbol E$向上，$\boldsymbol B$出纸面};
 \begin{scope}[xshift=7cm]
 \foreach \row/\sgn/\name in {0/1/{质子},1/-1/{电子}}{
  \begin{scope}[yshift={-3.4*\row cm}]
   \draw[DeepBlue,thick,domain=0:360,samples=200,variable=\t]
    plot({3-\sgn*.45*(pi*\t/180-sin(\t))},{-.45*(1-cos(\t))});
   \draw[->] ({3+\sgn*.7},-1.4)--({3-\sgn*.7},-1.4);
   \node at(3,-1.8){$\boldsymbol v_g$};
  \end{scope}}
 \node at(3,2){$\boldsymbol g$向下，$\boldsymbol B$出纸面};
 \end{scope}
\end{tikzpicture}
\par\smallskip{\small 各图分别归一化长度尺度，不比较回旋半径大小；电漂移同向，重力漂移随电荷改号而反向。}
\end{rTwoFigure}}

\renewcommand{\cyclotronDiagram}{\begin{rTwoFigure}
\begin{tikzpicture}[scale=.9,>=Latex]
 \filldraw[fill=PaleBlue,draw=DeepBlue] (-.12,-2) arc(-90:-270:2)--cycle;
 \filldraw[fill=SoftGray,draw=DeepBlue] (.12,-2) arc(-90:90:2)--cycle;
 \node at(-1.65,.9){D};\node at(1.65,.9){D};
 \node at(1.45,1.4){$\fieldOut\ B$};
 \def\rTwoY{0}
 \foreach \k in {1,...,8}{
  \pgfmathsetmacro{\rr}{.42*sqrt(\k)}
  \ifodd\k
   \draw[DeepBlue,thick] (-.12,\rTwoY)--(.12,\rTwoY);
   \draw[DeepBlue,thick] (.12,\rTwoY) arc(90:-90:\rr);
   \pgfmathparse{\rTwoY-2*\rr}\xdef\rTwoY{\pgfmathresult}
  \else
   \draw[DeepBlue,thick] (.12,\rTwoY)--(-.12,\rTwoY);
   \draw[DeepBlue,thick] (-.12,\rTwoY) arc(-90:-270:\rr);
   \pgfmathparse{\rTwoY+2*\rr}\xdef\rTwoY{\pgfmathresult}
  \fi
 }
 \fill (0,0) circle(2pt);
 \draw[->] (.12,.07)--(.43,-.10);
 \draw[thin] (-.12,-2)--(-.5,-2.55)--(-.5,-2.9);
 \draw[thin] (.12,-2)--(.5,-2.55)--(.5,-2.9);
 \draw (-.8,-2.9) rectangle(.8,-3.5);\node at(0,-3.2){高频振荡器};
 \node[anchor=west,align=left] at(2.7,.6){盒内：定半径圆弧\\[3pt]过缝：动能增加 $|q|\Delta U$\\[3pt]$r_k\propto\sqrt{k}$};
 \node[anchor=west,align=left] at(2.7,-1.9){窄缝用小间隔表示；\\实际加速轨迹连续。};
 \draw[thin] (0,2.15)--(0,2.5) node[above]{加速缝};
\end{tikzpicture}
\par\smallskip{\small 俯视图中磁场出纸面，正粒子顺时针弯转；半径只在过缝加速后增大。}
\end{rTwoFigure}}

\renewcommand{\splitProductsDiagram}{\begin{rTwoFigure}
\begin{tikzpicture}[scale=.9,>=Latex]
\begin{scope}[xshift=-3.6cm]
 \draw[DeepBlue,thick] (0,0) arc(90:330:1.3);
 \draw[Ink,dashed,thick] (0,0) arc(90:-30:1.3);
 \coordinate(P) at(0,0);\coordinate(C) at(0,-1.3);
 \coordinate(Q) at({sqrt(3)*1.3/2},{-1.5*1.3});
 \draw[guide] (P)--(C)--(Q);
 \fill(P) circle(2pt) node[above=12pt]{$P_0$};
 \fill(Q) circle(2pt) node[right]{$P_1$};
 \fill(C) circle(2pt) node[left]{$C_0$};
 \draw[vector] (P)--(-1.2,0) node[above]{$2v_0$};
 \draw[vector] (P)--(.6,0) node[above right]{$v_0$};
 \node at(-1.8,-1.2){$\fieldInto\ B$};
 \node[align=center] at(0,-3.5){实线：$+q,m$；虚线：$-q,2m$\\首段同圆反向，$\Delta t=2T/3$};
\end{scope}
\begin{scope}[xshift=3.1cm,yshift=-1.3cm,scale=.9]
 \foreach \k in {0,...,5}{
  \coordinate(P\k) at({1.3*cos(180+60*\k)},{1.3*sin(180+60*\k)});
  \fill(P\k) circle(2pt);
  \node at({1.65*cos(180+60*\k)},{1.65*sin(180+60*\k)}){$P_{\k}$};
 }
 \draw[DeepBlue,->] (P0)--(P1);\draw[guide] (P1)--(P2)--(P3)--(P4)--(P5)--(P0);
 \node[align=center] at(0,-2.5){六个碰撞点（连线不是轨迹）\\$P_6=P_0$，总周期 $4T$};
\end{scope}
\end{tikzpicture}
\end{rTwoFigure}}

\renewcommand{\beamFocusDiagram}{\begin{rTwoFigure}
\begin{tikzpicture}[x=.8cm,y=.8cm,>=Latex]
 \foreach \k/\phase/\lab in {0/270/{未聚焦},1/360/{首个近轴焦点}}{
 \begin{scope}[yshift={-3.8*\k cm}]
 \draw[axis] (0,0)--(8.8,0) node[right]{$z$};
 \draw[thick] (8,-1.4)--(8,1.4) node[above]{屏};
 \foreach \amp in {-.48,-.24,.24,.48}{
  \draw[DeepBlue,domain=0:8,samples=150,variable=\t] plot(\t,{\amp*(1-cos(\phase*\t/8))});
 }
 \fill(0,0) circle(2pt) node[below left]{$O$};
 \draw[<->] (0,-1.3)--(8,-1.3) node[midway,below]{$L$};
 \node[anchor=west] at(9,.6){\lab};
 \ifnum\k=0
 \draw[<->] (8.35,-.48)--(8.35,.48);
 \node[anchor=west] at(9,-.3){投影包络};
 \fi
 \end{scope}}
\end{tikzpicture}
\par\smallskip{\small 同一组近轴螺旋束在含轴平面上的投影；屏上最大径向位移需由两个横向分量合成。}
\end{rTwoFigure}}

\renewcommand{\magneticBilliardDiagram}{\begin{rTwoFigure}
\begin{tikzpicture}[scale=1.05,>=Latex]
 \def\R{2.3}\def\rc{1}\def\ang{30}
 \pgfmathsetmacro{\del}{2*atan2(cos(\ang),\R+sin(\ang))}
 \coordinate(O) at(0,0);\coordinate(P) at(0,-\R);
 \coordinate(C) at({-cos(\ang)},{-\R-sin(\ang)});
 \coordinate(Q) at({-\R*sin(\del)},{-\R*cos(\del)});
 \draw[curve] (O) circle(\R);
 \draw[DeepBlue,very thick] (P) arc(30:{150-\del}:1);
 \draw[guide] (P)--(O)--(Q) (C)--(P) (C)--(Q);
 \draw[vector] (P)--++(-.5,.866) node[above left]{$\boldsymbol v$};
 \fill(P) circle(2pt) node[below right]{$P$（小孔）};
 \fill(Q) circle(2pt) node[left]{$Q$};\fill(C) circle(2pt) node[below left]{$C$};
 \fill(O) circle(2pt) node[above]{$O$};
 \draw (0,-1.86) arc(90:120:.44) node[midway,above]{$\theta$};
 \draw (0,-.65) arc(-90:{-90-\del}:.65);
 \node at(-.26,-.89){$\delta$};
 \node[anchor=west,align=left] at(2.8,.4){$OP=OQ=R$\\$CP=CQ=r_{\rm c}$\\$CP\perp\boldsymbol v$\\[4pt]$\fieldInto\ B,\quad q>0$};
\end{tikzpicture}
\end{rTwoFigure}}

\renewcommand{\crossedFieldCycloidDiagram}{\begin{rTwoFigure}
\begin{tikzpicture}[x=.85cm,y=.85cm,>=Latex]
 \draw[thick] (-.3,2.3)--(9.5,2.3) (-.3,-.1)--(9.5,-.1);
 \foreach \x in {0,1,...,9}{\node at(\x,2.65){$+$};\node at(\x,-.4){$-$};}
 \draw[DeepBlue,very thick,domain=0:720,samples=300,variable=\t]
 plot({.72*(pi*\t/180-sin(\t))},{2.3-.72*(1-cos(\t))});
 \fill(0,2.3) circle(2pt);\node[above left] at(0,2.8){$v(0)=0$};
 \node at(4.5,1.5){$\fieldInto\ B$};
 \draw[<->] (9.3,2.3)--(9.3,.86) node[midway,right]{$h_{\max}$};
\end{tikzpicture}
\par\smallskip{\small 正电荷从上板静止释放；每个尖点处速度为零，$h_{\max}=2mE/(qB^2)$。}
\end{rTwoFigure}}

\renewcommand{\gravityDriftDiagram}{\begin{rTwoFigure}
\begin{tikzpicture}[x=1cm,y=1cm,>=Latex]
 \draw[thick] (-.2,0)--(8.8,0) node[right]{地面};
 \draw[DeepBlue,thick,domain=0:720,samples=300,variable=\t]
 plot({7+.3*(sin(\t)-pi*\t/180)+(cos(\t)-1)},{2+.3*(cos(\t)-1)-sin(\t)});
 \fill(7,2) circle(2pt) node[above right]{初始位置};
 \draw[vector] (7,2)--(7,1.3) node[right]{$v_0$};
 \draw[vector] (5.8,2.5)--(4.7,2.5) node[midway,above]{$v_D$};
 \draw[<->] (8.2,0)--(8.2,2) node[midway,right]{$H$};
 \node at(2.5,3.15){$q>0,\quad\fieldOut\ B$};
\end{tikzpicture}
\par\smallskip{\small 本图取不触地的一组参数；开始时切线竖直向下，回旋中心向左漂移。}
\end{rTwoFigure}}

\renewcommand{\alternatingLayersDiagram}{\begin{rTwoFigure}
\begin{tikzpicture}[x=.68cm,y=.68cm,>=Latex]
 \foreach \k in {0,2,4}{\fill[PaleBlue] (-1,\k) rectangle(8,\k+1);}
 \foreach \k in {1,3,5}{\fill[SoftGray] (-1,\k) rectangle(8,\k+1);}
 \foreach \k in {0,...,6}{\draw[RuleGray] (-1,\k)--(8,\k);}
 \draw[DeepBlue,very thick] plot file {\rTwoDataPath/layers.dat};
 \foreach \y in {.5,2.5,4.5}{\draw[vector](7,\y-.2)--(7,\y+.2);\node[right] at(7.2,\y){$E$};}
 \foreach \y/\n in {1.5/1,3.5/2,5.5/3}{\node at(7,\y){$\fieldInto$};\node[left]at(-1,\y){第\n 磁场层};}
 \draw[<->] (8.5,0)--(8.5,1) node[midway,right]{$d$};
\end{tikzpicture}
\par\smallskip{\small 图取 $\Lambda=2.5$，在第三磁场层内转向；水平速度不会在界面处突变。}
\end{rTwoFigure}}

\renewcommand{\dampedSpiralDiagram}{\begin{rTwoFigure}
\begin{tikzpicture}[scale=1.05,>=Latex]
 \draw[axis] (-1,0)--(3,0) node[right]{$x$};\draw[axis] (0,-.3)--(0,3.7) node[above]{$y$};
 \draw[DeepBlue,thick,domain=0:1260,samples=450,variable=\t]
 plot({2/(1+.22^2)*(.22+exp(-.22*pi*\t/180)*(-.22*cos(\t)+sin(\t)))},
 {2/(1+.22^2)*(1+exp(-.22*pi*\t/180)*(-.22*sin(\t)-cos(\t)))});
 \fill(0,0) circle(2pt) node[below left]{$O$};
 \fill({.44/(1+.22^2)},{2/(1+.22^2)}) circle(2pt);
 \draw[thin] ({.44/(1+.22^2)},{2/(1+.22^2)})--(2.8,2.1) node[right]{$\boldsymbol r_\infty$};
 \draw[vector] (0,0)--(1,0) node[below]{$v_0$};
 \node[anchor=west,align=left] at(3.8,1){$\fieldInto\ B,\quad q>0$\\初始速度沿 $+x$\\终点有 $y_\infty>0$};
\end{tikzpicture}
\end{rTwoFigure}}

\renewcommand{\dampedDriftDiagram}{\begin{rTwoFigure}
\begin{tikzpicture}[x=.7cm,y=.7cm,>=Latex]
 \pgfmathsetmacro{\ux}{1/(1+.25^2)}\pgfmathsetmacro{\uy}{.25/(1+.25^2)}
 \pgfmathsetmacro{\cr}{((3-\ux)*.25-\uy)/(1+.25^2)}
 \pgfmathsetmacro{\ci}{(-(3-\ux)-.25*\uy)/(1+.25^2)}
 \draw[axis] (-.3,0)--(12.6,0) node[right]{$x$};\draw[axis](0,-3)--(0,2)node[above]{$y$};
 \draw[guide] (0,{\ci-\uy*\cr/\ux})--(12,{\ci+\uy*(12-\cr)/\ux});
 \draw[DeepBlue,thick,domain=0:720,samples=350,variable=\t]
 plot({\ux*pi*\t/180+\cr*(1-exp(-.25*pi*\t/180)*cos(\t))-\ci*exp(-.25*pi*\t/180)*sin(\t)},
 {\uy*pi*\t/180+\ci*(1-exp(-.25*pi*\t/180)*cos(\t))+\cr*exp(-.25*pi*\t/180)*sin(\t)});
 \draw[vector](0,0)--(1,0)node[above]{$v_0$};
 \node[anchor=west]at(6,-2.5){$E\parallel+y,\quad\fieldOut\ B$};
 \draw[vector](9,.8)--(10.5,1.175)node[above]{$\boldsymbol v_\infty$};
\end{tikzpicture}
\par\smallskip{\small 虚线是长期直线渐近线；两个方向的回旋瞬态均指数衰减。}
\end{rTwoFigure}}
